\BourbakiExerciseGroup

\begin{enumerate}
\def\labelenumi{\arabic{enumi})}
\item
  Show that if A is an integral domain then the polynomial ring \(A[X_{\iota}]_{\iota \in I}\) is not a principal ideal domain when one of the following conditions holds: 1) \(\operatorname{Card}(I) \geqslant 2\) ; 2) A is not a field (consider the ideals \((\mathbf{X}_{\alpha}) + (\mathbf{X}_{\beta})\) for \(\alpha \neq \beta\) , and \((a) + (\mathbf{X}_{\alpha})\) where \(a \neq 0\) is a noninvertible element of A).
\item
  Show that, in the ring \(\mathbf{K}[[\mathbf{X}]]\) of formal power series over a field \(\mathbf{K}\) , every irreducible element is an associate of \(\mathbf{X}\) .
\item
  Let \(\mathbf{A}\) be an integral domain in which every nonempty set of ideals has a maximal element and in which every maximal ideal is principal; show that \(\mathbf{A}\) is a principal ideal domain. (Notice that if \(\mathfrak{a} \neq \mathbf{A}\) is an ideal of \(\mathbf{A}\) then there exists a maximal ideal \((p) \supset \mathfrak{a}\) such that \(\frac{1}{p} \mathfrak{a}\) , which is an ideal of \(\mathbf{A}\) , strictly contains \(\mathfrak{a}\) .)
\end{enumerate}

\begin{enumerate}
\def\labelenumi{\alph{enumi})}
\setcounter{enumi}{1}
\tightlist
\item
  Let K be a field, let \(\mathrm{K}_{1}=\mathrm{K}((\mathrm{X}))\) be the field of formal power series in one indeterminate X over K (IV, p.~38), and let \(B=K_{1}[Y]\) be the ring of formal power series in Y with coefficients in \(K_{1}\) . Let A be the subring of B consisting of those power series \(\sum_{n=0}^{\infty}a_{n}(\mathrm{X})\mathrm{Y}^{n}\) such that the power series \(a_{0}(\mathrm{X})\) involves only powers of X with exponent \(\geqslant 0\) . Show that the principal ideal (X) is the unique maximal ideal of A, but the ideal of A generated by the elements \(YX^{-n}\) ( \(n \geqslant 0\) an arbitrary integer) is not principal.
\end{enumerate}

\begin{enumerate}
\def\labelenumi{\arabic{enumi})}
\setcounter{enumi}{3}
\tightlist
\item
  Let A be an integral domain and let S be a multiplicatively closed subset of A not containing 0. Let \(S^{-1}A\) denote the ring of elements \(s^{-1}a\) ( \(s \in S\) , \(a \in A\) ) in the field of fractions K of A (I, p.~116).
\end{enumerate}

\begin{enumerate}
\def\labelenumi{\alph{enumi})}
\item
  Show that if \(\mathbf{A}\) is a principal ideal domain then so is \(\mathbf{S}^{-1}\mathbf{A}\) .
\item
  Let A be a principal ideal domain and let p be an irreducible element of A; show that the complement S of the ideal \((p)\) is multiplicatively closed; in the ring \(S^{-1}A\) the ideal \(pS^{-1}A\) is the unique maximal ideal, and the quotient field \(S^{-1}A/pS^{-1}A\) is isomorphic to \(A/(p)\) .
\end{enumerate}

\begin{enumerate}
\def\labelenumi{\arabic{enumi})}
\setcounter{enumi}{4}
\tightlist
\item
  Let \(\mathbf{A}\) be a commutative ring in which every ideal is finitely generated.
\end{enumerate}

\begin{enumerate}
\def\labelenumi{\alph{enumi})}
\item
  An element \(p \neq 0\) of \(\mathbf{A}\) is said to be indivisible if it is not invertible and if the relation \(p = xy (x, y \in \mathbf{A})\) implies that either \(x\) or \(y\) belongs to \(A p\) . Show that every element \(a \neq 0\) of \(\mathbf{A}\) can be written (in at least one way) as a product of indivisible elements and of one invertible element (use Lemma 1 of VII, p.~4).
\item
  An idempotent e in a commutative ring A is said to be indecomposable if there does not exist a pair of nonzero idempotents f, g such that \(f + g = e\) and fg = 0. Two distinct indecomposable idempotents have product zero. Show that every idempotent in the ring A is a sum of indecomposable idempotents (show that if \(e\) is an idempotent then so is \(1 - e\) , and use Lemma 1 of VII, p.~4). Deduce that \(\mathbf{A}\) is the direct product of a finite family of rings \(\mathbf{A}_i\) , in each of which every ideal is finitely generated, and the unit element is the only nonzero idempotent. Suppose this family has at least two elements. Then an element \(a = \sum_{i} a_i\) (with \(a_i \in \mathbf{A}_i\) for all \(i\) ) is indivisible if and only if there is an index \(k\) such
\end{enumerate}

that \(a_{k}\) is either indivisible or zero in \(\mathbf{A}_{k}\) , and \(a_{i}\) is invertible in \(\mathbf{A}_{i}\) for all \(i \neq k\) ; if \(a_{k} = 0\) then \(\mathbf{A}_{k}\) is necessarily an integral domain.

\begin{enumerate}
\def\labelenumi{\alph{enumi})}
\setcounter{enumi}{2}
\item
  Show that if 1 is the only nonzero idempotent of \(\mathbf{A}\) , and if \(x \in \mathbf{A}\) is such that \(x^k \in \mathbf{A}x^{k+1}\) for some integer \(k \geqslant 1\) , then either \(x\) is invertible or \(x^k = 0\) (if \(x^k = ax^{k+1}\) , then consider \(a^k x^k\) ).
\item
  Let \(p \neq 0\) be a zero divisor in A. Show that there exist a finite sequence \((a_k)_{1 \leqslant k \leqslant m}\) of nonzero elements of A such that \(0 = pa_1\) , \(a_k = pa_{k+1}\) for \(1 \leqslant k < m\) , \(a_m \notin \mathbf{A}_p\) , and \((a_k) \neq (a_{k+1})\) for \(1 \leqslant k < m\) .
\end{enumerate}

\begin{enumerate}
\def\labelenumi{\arabic{enumi})}
\setcounter{enumi}{5}
\tightlist
\item
  A commutative ring A is called a principal ideal ring if every ideal of A is principal. a) Show that a quotient ring of a principal ideal ring and a finite product of principal ideal rings are principal ideal rings.
\end{enumerate}

\begin{enumerate}
\def\labelenumi{\alph{enumi})}
\setcounter{enumi}{1}
\item
  Show that every principal ideal ring A is a direct product of a finite family \((\mathbf{A}_{i})_{1\leq i\leq m}\) of principal ideal rings in each of which the unit element is the only nonzero idempotent (use Ex. 5, b)).
\item
  Suppose that A is a principal ideal ring in which the unit element is the only nonzero idempotent, and which contains zero divisors other than 0. Show that there exists an indivisible nilpotent element p in A. (Using Ex. 5, a), show that there exists an indivisible element \(p \neq 0\) in A which is a zero divisor; apply Ex. 5 d), then prove that the principal ideal Ab of \(y \in A\) such that \(yp^{m} \in Ap^{m+1}\) is the whole of A; finally apply Ex. 5, c)). d) With the hypotheses and notation as in c), suppose that \(p^{m} = 0\) and \(p^{m-1} \neq 0\) . Show that the annihilator of \(p^{m-1}\) is Ap (otherwise it would be a principal ideal Ac with \(p \in Ac\) and \(c \notin Ap\) ; using the fact that p is indivisible, deduce that c would be invertible); deduce that, for \(1 \leq k \leq m-1\) , the annihilator of \(p^{m-k}\) is \(Ap^{k}\) . Conclude that the ideal Ap is maximal (argue as for the first assertion of d), using the fact that 1 - xp is invertible in A for all \(x \in A\) .
\item
  Deduce from \(d\) ) that for all \(x \in \mathbf{A}\) there exists an integer \(k \geqslant 0\) and an invertible element \(u\) such that \(x = up^k\) , and that the only ideals of \(\mathbf{A}\) are the \(Ap^k\) ( \(0 \leqslant k \leqslant m\) ) (observe that every element of \(\mathbf{A}\) not belonging to \(Ap\) is invertible in \(\mathbf{A}\) ).
\item
  Let A be a principal ideal ring, let \((x_{i})\) be a family of elements of A, and let \(Ad\) be the ideal generated by the family \((x_{i})\) ; show that there exist elements \(x_{i}^{\prime}\) of A such that \(x_{i} = dx_{i}^{\prime}\) for all \(i\) , and \(\sum_{i} Ax_{i}^{\prime} = A\) (reduce to the case where 1 is the only
\end{enumerate}

nonzero idempotent of A).

\(g)\) Give an example of a principal ideal ring and an indivisible element \(p\) of \(\mathbf{A}\) such that the ideal \(Ap\) is not maximal (use Ex. 5, \(b\) )).

\begin{enumerate}
\def\labelenumi{\arabic{enumi})}
\setcounter{enumi}{6}
\tightlist
\item
  Given an integral domain \(\mathbf{A}\) , we call a map \(w\) from \(\mathbf{A}' = \mathbf{A} - \{0\}\) into \(\mathbf{N}\) a Euclidean function if it satisfies the following conditions:
\end{enumerate}

\((S_{1}) w(xy) \geqslant w(y)\) for each pair \(x, y\) of elements of \(\mathbf{A}'\) .

(S \(_{II}\) ) If \(a \in \mathbf{A}'\) and \(b \in \mathbf{A}'\) , then there exist elements \(q\) , \(r\) of \(\mathbf{A}\) such that \(a = bq + r\) , and such that either \(r = 0\) or \(w(r) < w(b)\) .

If there exists a Euclidean function on A then A is called a Euclidean domain.

\begin{enumerate}
\def\labelenumi{\alph{enumi})}
\item
  Show that \(\mathbf{Z}\) and \(\mathbf{K}[\mathbf{X}]\) (where \(\mathbf{K}\) is a field) are Euclidean domains.
\item
  Show that every Euclidean domain is a principal ideal domain (if \(\mathfrak{a}\) is an ideal of \(\mathbf{A}\) , choose \(a \in \mathfrak{a}\) such that \(w(a)\) is as small as possible).
\item
  Show that the ring A of Gaussian integers (VII, p.~1) is a Euclidean domain with respect to the function \(w(a + bi) = a^2 + b^2 = \mathbf{N}_{\mathbf{Q}(i) / \mathbf{Q}}(a + bi)\) (notice that, for all \(x \in \mathbf{Q}(i)\) , there exists \(y \in \mathbf{A}\) such that \(\mathbf{N}(x - y) \leqslant 1/2\) ).
\end{enumerate}

\(d\) ) Show that the quadratic extension ring \(\mathbf{B}\) of \(\mathbf{Z}\) with basis \((1,e)(e^2 = 2)\) , is a Euclidean domain with respect to the function

\[
w (a + b e) = \left| a ^ {2} - 2 b ^ {2} \right| = \left| N _ {Q (e) / Q} (a + b e) \right|
\]

(for all \(x \in \mathbf{Q}(e)\) , there exists \(y \in B\) such that \(|\mathbf{N}(x - y)| \leqslant 1/2\) ).

\begin{enumerate}
\def\labelenumi{\arabic{enumi})}
\setcounter{enumi}{7}
\tightlist
\item
  Let \(\mathbf{K}\) be a quadratic field, that is an extension of \(\mathbf{Q}\) of degree 2.
\end{enumerate}

\begin{enumerate}
\def\labelenumi{\alph{enumi})}
\item
  Show that \(\mathbf{K} = \mathbf{Q}(e)\) where \(e^2\) is a rational integer \(d \neq 1\) not divisible by a square \(>1\) (in \(\mathbf{Z}\) ).
\item
  If \(\alpha = a + be\) ( \(a, b \in \mathbf{Q}\) ), let \(\overline{\alpha}\) denote the conjugate \(a - be\) of \(\alpha\) . We say that \(\alpha\) is an integer of the quadratic field \(\mathbf{K}\) if its norm \(\alpha\overline{\alpha}\) and its trace \(\alpha + \overline{\alpha}\) are rational integers. Show that the integers of \(\mathbf{K}\) form a ring \(\mathbf{A}\) (notice that if \(\alpha\) and \(\beta\) are integers of \(\mathbf{K}\) then the sum and product of the rational numbers \(\alpha\beta + \overline{\alpha\beta}\) and \(\alpha\overline{\beta} + \overline{\alpha}\beta\) are integers, hence that these rational numbers are the roots of a monic polynomial over \(\mathbf{Z}\) ; deduce that they are rational integers).
\item
  Show that the integers of \(\mathbf{K} = \mathbf{Q}(e)\) form a \(\mathbf{Z}\) -module with basis \((1, e)\) if \(d - 1 \not\equiv 0 \pmod{4}\) , or \((1/2(1 + e), 1/2(1 - e))\) if \(d - 1 \equiv 0 \pmod{4}\) (if \(a + be\) is an integer of \(\mathbf{K}\) , check that \(2a\) and \(2b\) are rational integers such that \((2a)^2 - d(2b)^2 \equiv 0(4)\) , and investigate whether they can be odd). Show that the discriminants of these bases (III, p.~549, Def. 3) are \(4d\) and \(d\) respectively.
\end{enumerate}

\begin{enumerate}
\def\labelenumi{\arabic{enumi})}
\setcounter{enumi}{8}
\item
  * (Minkowski's Theorem) If \(a, b, c, d\) are real numbers, and \(\mathbf{D} = \begin{vmatrix} a & b \\ c & d \end{vmatrix}\) , then the inequalities \(|na + mb| \leqslant \mathbf{A}\) , \(|nc + md| \leqslant \mathbf{B}\) have nonzero integer solutions \((n, m)\) if \(\mathbf{A} > 0\) , \(\mathbf{B} > 0\) , \(\mathbf{D} > 0\) and \(\mathbf{AB} \geqslant \mathbf{D}\) , and such solutions are finite in number (consider the discrete subgroup \(\mathbf{G}\) of \(\mathbf{R}^2\) generated by \((a, c)\) and \((b, d)\) , and show that the rectangle \(\mathbf{E}\) defined by the relations \(0 \leqslant x \leqslant \mathbf{A}\) , \(0 \leqslant y \leqslant \mathbf{B}\) , whose area is greater than that of the parallelogram constructed on the vectors \((a, c)\) and \((b, d)\) , contains two distinct vectors which are congruent mod \(\mathbf{G}\) ; argue by contradiction, noticing that otherwise the \(p^2\) rectangles obtained from \(\mathbf{E}\) by means of the translations ( \(ha + kb\) , \(hc + kd\) ), where \(0 \leqslant h < p\) and \(0 \leqslant k < p\) , would be pairwise disjoint).
\item
  Let \(d\) be a rational integer not divisible by a square \(>1\) . In the quadratic field \(\mathbf{K} = \mathbf{Q}(\sqrt{d})\) , we shall examine the units, that is to say the invertible elements of the ring \(\mathbf{A}\) of integers of \(\mathbf{K}\) (Ex. 8, b)).
\end{enumerate}

\begin{enumerate}
\def\labelenumi{\alph{enumi})}
\item
  If \(d < 0\) (« imaginary quadratic fields »), then show that the only units of \(K\) are 1 and -1, except for \(d = -1\) , when \(i\) and -i are also units, and for \(d = -3\) , where the four numbers \(1/2(\pm 1 \pm \sqrt{-3})\) are also units.
\item
  Suppose from now on that \(d > 0\) (« real quadratic fields »), and suppose that \(\mathbf{K}\) is embedded in a maximal ordered extension of \(\mathbf{Q}\) (VI, p.~25), * for example in \(\mathbf{R}\) . * Show that the positive units form a subgroup \(U'\) of the group U of units of K, and that U is the direct product of \(U'\) and \(\{-1,1\}\) .
\item
  Show that the units of K are the integers with norm 1 or -1.
\end{enumerate}

\(d\) ) Show that the group \(\mathbf{U}^{\prime}\) is nontrivial (let D denote the discriminant of the integer basis (Ex. 8, \(c\) ) ; show that the integers \(\alpha\) of K such that \(\mathbf{N}(\alpha)\leqslant \mathbf{D}\) are partitioned into a finite number of congruence classes mod U ; choose an integer \(\beta_{i}\) in each class ; if B is a rational number such that \(0 < \mathrm{B} < \inf (|\beta_i|)\) , show (Ex. 9) that there exists an integer \(\beta\) of K such that \(0 < |\beta |\leqslant \mathrm{B}\) and \(\mathbf{N}(\beta)\leqslant \mathbf{D}\) ; deduce that some \(\beta \beta_{i}^{-1}\) is a unit \(\eta\) such that \(|\eta| < 1\) ).

\begin{enumerate}
\def\labelenumi{\alph{enumi})}
\setcounter{enumi}{4}
\item
  Show that \(\mathbf{U}'\) is isomorphic to \(\mathbf{Z}\) (show that there exists a unit \(\varepsilon\) such that \(|\varepsilon| < 1\) and \(|\varepsilon|\) is as great as possible with this property, using Ex. 9 and the fact that \(|\overline{\xi}| = |\xi|^{-1}\) for a unit).
\item
  Show that in \(\mathbf{Q}(\sqrt{2})\) , the element \(\sqrt{2} - 1\) is a generator of \(\mathbf{U}'\) , and has norm \(-1\) .
\end{enumerate}

\(g\) ) Show that if \(d\) has a prime factor \(p\) of the form \(4n - 1\) , then every generator of \(\mathbf{U}'\) has norm 1. (Observe that \(-1\) cannot be a square in the field \(\mathbf{F}_p\) , using V, p.~93, Prop. 1, c)).

¶ 11) We investigate the quadratic fields \(\mathbf{Q}(\sqrt{d})\) for which the map \(\alpha \mapsto |\mathrm{N}(\alpha)|\) is a Euclidean function on the integers of the field (Ex. 7).

\begin{enumerate}
\def\labelenumi{\alph{enumi})}
\item
  Show that a necessary and sufficient condition for this is that, for all \(\alpha \in \mathbf{Q}(\sqrt{d})\) , there exist an integer \(\beta\) of the field such that \(|\mathrm{N}(\alpha - \beta)| < 1\) ; write down this condition, expressing \(\alpha\) and \(\beta\) in terms of an integer basis of the field (Ex. 8, c)).
\item
  In the case of an imaginary quadratic field \((d < 0)\) , show that the only values of d for which \(|\mathbf{N}(\alpha)|\) is a Euclidean function are -1, -2, -3, -7, -11.
\item
  For real quadratic fields, we will restrict ourselves to the cases \(d \equiv 2\) or 3 (mod 4). Show that there are then only finitely many values of \(d\) such that \(|\mathrm{N}(\alpha)|\) is a Euclidean function. (If \(|\mathrm{N}(\alpha)|\) is a Euclidean function and \(a\) is a natural number such that \(0 \leqslant a \leqslant d\) and \(a \equiv t^2(d)\) , then show that either \(a\) or \(d - a\) has the form \(x^2 - dy^2\) , where \(x\) and \(y\) are rational integers, by considering the element \(t / \sqrt{d}\) of \(\mathbf{Q}(\sqrt{d})\) ; next show that for large enough \(d\) there exists an odd integer \(z\) such that \(5d < z^2 < 6d\) (if \(d \equiv 3 \mod 4\) ) or \(2d < z^2 < 3d\) (if \(d \equiv 2 \mod 4\) ); then take \(a = z^2 - d[z^2/d]\) and show that, under these conditions, neither \(a\) nor \(d - a\) has the form \(x^2 - dy^2\) , by examining residues mod 8). Show that the only values of \(d\) for which no such integer \(z\) exists are 3, 7, 11, 19, 35, 47 and 59 (for \(d \equiv 3 \mod 4\) ) and 2, 6, 14 and 26 (for \(d \equiv 2 \mod 4\) ). Show that there is indeed a Euclidean function for \(d = 2\) and \(d = 3\) .
\end{enumerate}

\begin{enumerate}
\def\labelenumi{\arabic{enumi})}
\setcounter{enumi}{11}
\tightlist
\item
  Show that the ring of integers of the quadratic field \(\mathbf{Q}(\alpha)\) , where \(\alpha^2 = -5\) , is not a principal ideal domain (see VI, p.~34, Ex. 27). Same question for \(\mathbf{Q}(\beta)\) , where \(\beta^2 = 10\) (we have
\end{enumerate}

\[
2. 3 = (4 + \beta) (4 - \beta)).
\]

\begin{enumerate}
\def\labelenumi{\arabic{enumi})}
\setcounter{enumi}{12}
\item
  Let \(f\) be a nonconstant polynomial in one indeterminate over \(\mathbf{Z}\) . Show that for each integer \(k \geqslant 1\) there exist infinitely many integers \(n \in \mathbf{N}\) such that \(f(n)\) is nonzero and divisible by at least \(k\) distinct primes (one can proceed by induction on \(k\) : if \(n \in \mathbf{N}\) is such that \(f(n)\) is nonzero and divisible by at least \(k\) distinct primes, consider the polynomial \(f(n + Xf(n)^2)\) .
\item
  The set of prime numbers of the form \(4n - 1\) (resp. \(6n - 1\) ) is infinite (same method as Prop. 5 (VII, p.~5); notice that every prime number except 2 (resp. except 2 and 3) has the form \(4n \pm 1\) (resp. \(6n \pm 1\) )).
\item
  A number of the form \(2^k + 1\) ( \(k \geqslant 1\) an integer) can be a prime only if \(k\) is a power of 2. If a prime number \(p\) divides \(2^{2^n} + 1\) , then it cannot divide \(2^{2^m} + 1\) for \(m > n\) ( \(2^{2^m} \equiv 1(p)\) ); use this to give a new proof of Prop. 5. Show that \(2^{2^5} + 1\) is not prime (put \(a = 2^7\) and \(b = 5\) ; show that \(1 + ab - b^4 = 2^4\) and that \(1 + ab\) divides
\end{enumerate}

\[
n = 2 ^ {4} a ^ {4} + 1 = (1 + a b - b ^ {4}) a ^ {4} + 1.
\]

\begin{enumerate}
\def\labelenumi{\arabic{enumi})}
\setcounter{enumi}{15}
\item
  A number of the form \(a^k - 1\) ( \(a\) and \(k\) integers, \(a > 0\) , \(k > 1\) ) can be prime only if \(a = 2\) and \(k\) is prime.
\item
  Show that if \(a^n + 1\) is prime, where \(a > 1\) and \(n > 1\) , then \(a\) is even and \(n = 2^m\) .
\item
  Let M be the multiplicative monoid generated by a finite number of integers \(q_{i} > 1\) ( \(1 \leqslant i \leqslant m\) ). Show that, for each integer k \textgreater{} 0, there exist k consecutive integers which do not belong to M (if \(h_{1}, \ldots, h_{m}\) are m integers \textgreater{} 1, and k is the greatest number of consecutive integers none of which is divisible by any of the \(h_{i}\) , then \((k + 1)^{-1} \leqslant h_{1}^{-1} + \cdots + h_{m}^{-1}\) ; now notice that, for all r \textgreater{} 0, every sufficiently large number in M is divisible by at least one \(q_{i}^{+}\) ). Obtain a new proof of Prop. 5 of VII, p.~5 from this result.
\item
  For every integer \(n > 0\) , show that the exponent of a prime \(p\) in the prime factorisation of \(n!\) is \(\sum_{k=1}^{\infty} [np^{-k}]\) (which has only finitely many nonzero terms).
\item
  * Let \(\pi(x)\) be the number of primes less than the real number \(x > 0\) . For each integer \(n > 0\) , show that
\end{enumerate}

\[
n ^ {\pi (2 n) - \pi (n)} \leqslant \binom {2 n} {n} \leqslant (2 n) ^ {\pi (2 n)}
\]

(notice that \(\binom{2n}{n}\) is a multiple of the product of all the primes \(q\) such that \(n < q \leqslant 2n\) ); on the other hand, use Ex. 19 to show that \(\binom{2n}{n}\) divides the product \(\prod p^{r(p)}\) , where \(p\) runs through the set of prime numbers \(\leqslant 2n\) , and where \(r(p)\) is the largest integer such that \(p^{r(p)} \leqslant 2n\) . Deduce from this result that there exist two real numbers \(a\) and \(b\) such that \(0 < a < b\) and

\[
a n (\log n) ^ {- 1} \leqslant \pi (n) \leqslant b n (\log n) ^ {- 1}
\]

(notice that \(2^{2n}(2n + 1)^{-1} \leqslant (2n)! (n!)^{-2} \leqslant 2^{2n}\) ). *

\(^{1}\) It has been shown that \(\pi(x) \sim \frac{x}{\log x}\) as x tends to \(+\infty\) (cf.~for example A. E. INGHAM, The distribution of prime numbers (Cambridge tracts, No.~30), Cambridge University Press, 1932).

\begin{enumerate}
\def\labelenumi{\arabic{enumi})}
\setcounter{enumi}{20}
\tightlist
\item
  Show that for \(m \geqslant 1\) and \(n \geqslant 1\) the rational number \(\sum_{k=n}^{n+m} k^{-1}\) is never an integer (if
\end{enumerate}

\(2^{q}\) is the largest power of 2 which divides at least one of the numbers \(n, n+1, \ldots, n+m\) , then it divides only one of these numbers).

\begin{enumerate}
\def\labelenumi{\arabic{enumi})}
\setcounter{enumi}{21}
\item
  Let \(a\) and \(b\) be two integers \(> 0\) , let \(d\) be the gcd of \(a\) and \(b\) , and let \(n\) be an integer \(\geqslant 1\) ; put \(a = da_1\) and \(b = db_1\) ; show that the denominator in the irreducible expression for the fraction \(a(a + b) \ldots (a + (n - 1)b) / n!\) admits only prime factors which divide \(b_1\) and are \(\leqslant n\) (observe that if \(p\) is a prime number not dividing \(b_1\) , then \(b_1\) is invertible mod \(p^r\) for all \(r > 0\) ). Hence obtain a direct proof of the fact that the binomial coefficients \(n! / (m!(n - m)! )\) are integers.
\item
  Let \(n > 1\) be a rational integer; show that the gcd of the \(n - 1\) binomial coefficients \(\binom{n}{k}\) ( \(1 \leqslant k \leqslant n - 1\) ) is equal to 1 if \(n\) is divisible by two distinct prime numbers, and is equal to \(p\) if \(n\) is a power of the prime number \(p\) .
\item
  \begin{enumerate}
  \def\labelenumii{\alph{enumii})}
  \tightlist
  \item
    Let \(n = \prod_{k} p_k^{\nu_k} > 0\) be a rational integer expressed in terms of its prime factors.
  \end{enumerate}
\end{enumerate}

Show that the sum of the positive divisors of n is given by the formula

\[
\sigma (n) = \prod_ {k} \frac {p _ {k} ^ {\nu_ {k} + 1} - 1}{p _ {k} - 1}.
\]

\begin{enumerate}
\def\labelenumi{\alph{enumi})}
\setcounter{enumi}{1}
\tightlist
\item
  We say that \(n\) is perfect if \(2n = \sigma(n)\) . Show that an even number \(n\) is perfect if and only if it has the form \(2^h(2^{h+1} - 1)\) , where \(h \geqslant 1\) and \(2^{h+1} - 1\) is prime. (If \(n\) has the form \(2^m q\) , where \(m \geqslant 1\) and \(q\) is odd, show that \(q\) is divisible by \(2^{m+1} - 1\) , and deduce that, if \(q\) were not prime, then we would have \(\sigma(n) > 2n\) ).
\end{enumerate}

\begin{enumerate}
\def\labelenumi{\arabic{enumi})}
\setcounter{enumi}{24}
\item
  Let \(a\) and \(b\) be two coprime integers \(> 0\) . If \(a\) and \(b\) are \(> 2\) , then there exist integers \(x\) and \(y\) such that \(ax + by = 1\) and \(|x| < 1/2b\) , \(|y| < 1/2a\) ; moreover the integers \(x\) and \(y\) are uniquely determined by these conditions. What happens in the case where one of \(a, b\) is equal to 2?
\item
  Let \(x, y\) be two coprime integers; then every integer \(z\) such that \(|z| < |xy|\) can be expressed, in either one or two ways, in the form \(z = ux + vy\) where \(u\) and \(v\) are integers with \(|u| < |y|\) and \(|v| < |x|\) . If \(z\) can be expressed in two different ways in this form, then neither \(x\) nor \(y\) can divide \(z\) ; the example \(x = 5, y = 7, z = 12\) shows that the converse fails.
\item
  Let \(f\) and \(g\) be two coprime polynomials in \(\mathbf{K}[\mathbf{X}]\) ; show that every polynomial \(h\) such that \(\deg(h) < \deg(fg)\) can be expressed uniquely in the form \(h = uf + vg\) where \(u\) and \(v\) are polynomials such that \(\deg(u) < \deg(g)\) and \(\deg(v) < \deg(f)\) . First give a proof analogous to that of Ex. 26 (based on the Euclidean algorithm), then give a proof in which the problem is considered as a linear problem.
\item
  \begin{enumerate}
  \def\labelenumii{\alph{enumii})}
  \tightlist
  \item
    Let K be an algebraically closed commutative field, let f be a nonconstant monic polynomial in K{[}X{]}, and let g be a polynomial in K{[}X, Y{]}; suppose that for each of the roots \(\alpha_{i}\) of f there exists a root \(\beta_{i}\) of \(g(\alpha_{i}, \mathbf{Y}) = 0\) such that \(\frac{\partial g}{\partial \mathbf{Y}} (\alpha_{i}, \beta_{i}) \neq 0\) . Then show that there exists a polynomial \(h(\mathbf{X}) \in \mathbf{K}[\mathbf{X}]\) such that \(g(\mathbf{X}, h(\mathbf{X})) \equiv 0 \pmod{f(\mathbf{X})}\) . (Using Prop. 4, reduce to the case \(f(\mathbf{X}) = \mathbf{X}^{m}\) (for some \(m \geqslant 1\) ); notice that the ring \(K[X]/(X^{m})\) is isomorphic to the quotient ring \(K[[X]]/(X^{m})\) , and use the Cor. of IV, p.~37).
  \end{enumerate}
\end{enumerate}

\begin{enumerate}
\def\labelenumi{\alph{enumi})}
\setcounter{enumi}{1}
\tightlist
\item
  In particular show that if m is an integer not divisible by the characteristic of K, then for every pair of coprime monic polynomials f and \(f_{0}\) in K{[}X{]}, there exists a polynomial \(h(\mathbf{X})\) in K{[}X{]} such that \((h(\mathbf{X}))^{m} \equiv f_{0}(\mathbf{X}) \pmod{f(\mathbf{X})}\) .
\end{enumerate}

